\begin{definition}[Zero-crossing physical graph]
    \label{def:peps_zero_crossing_graph}
    \lean{TNLean.PEPS.torusDualZeroCrossingGraph}\leanok
    Given an actual native dual path $p$, retain a horizontal physical bond
    based at $v$ exactly when its signed horizontal crossing count is zero,
    and an upward physical bond based at $v$ exactly when its signed downward
    crossing count is zero. Include both orientations of each retained bond.
    This gives a finite physical subgraph $Z_p$, independent of the flux group.
\end{definition}

\begin{lemma}[Zero-crossing edges are untwisted]
    \label{lem:peps_zero_crossing_transport}
    \lean{TNLean.PEPS.torusDualFlux_transport_eq_one_of_zero_crossing_edge}\leanok
    For any group and insertion element, the actual directed flux transport
    on every edge of $Z_p$ is identity.
\end{lemma}
\begin{proof}\leanok
    \uses{def:peps_zero_crossing_graph}
    The native transport is the insertion element to the signed crossing
    count, or its inverse according to orientation. A zero exponent gives
    identity in all four orientations, including periodic seams.
\end{proof}

\begin{lemma}[Vacuum density from zero-crossing boundary connectivity]
    \label{lem:peps_zero_crossing_density}
    \lean{TNLean.PEPS.IsGIsometric.torusFlux_normalized_density_eq_vacuum_of_zero_crossing_connectivity}
    \leanok
    Let a native torus with periods at least three carry a regular
    group-isometric tensor for a finite group. Suppose the induced graphs
    on a physical observation $R$ and its complement are connected and the
    boundary is nonempty. Assume the actual path has zero signed crossing
    counts on all bonds incident to $R$, and one component of the computed
    graph $Z_p$ in the complement contains every complementary boundary
    endpoint. Then both native coefficient cuts are nonzero and the normalized
    observed density equals vacuum. No simple connectedness is required.
\end{lemma}
\begin{proof}\leanok
    \uses{lem:peps_zero_crossing_transport,lem:peps_anchored_twist_density}
    Restrict $Z_p$ to the complement. Its edge transports are identity by the
    preceding lemma; the stated component condition places all complementary
    boundary endpoints in one untwisted component. Zero incident counts also
    give identity on every observed incident operator. Apply the noncentral
    complementary-contraction result and convert its ordered-edge cuts to
    the actual native cuts, including the vacuum cut.
\end{proof}
